% Generado automáticamente desde notas/03_sprint.md con md2tex.py; no editar a mano.
\begin{nota}
Estudio de viabilidad con pruebas reales (29 sprints con FastF1, 7 con OpenF1, Jolpica, FIA). Fecha: 2026-09-29. No se ha modificado código del repo; prototipos en el scratchpad.
\end{nota}

\section{Conclusiones}

\begin{enumerate}
\item \textbf{Viable y barato}: FastF1 (\code{get\_\allowbreak{}session(\allowbreak{}año,\allowbreak{} ronda,\allowbreak{} 'Sprint')}) tiene el 100 \% de las vueltas oficiales de los 29 sprints disputados (11 319 vueltas; 568/590 pilotos exactos y el resto son la vuelta de abandono, 1 DNS y 1 DSQ), con posición, compuesto, stint, boxes y estado de pista (SC/VSC/roja). Cubre los tres formatos (2021-22, 2023 Shootout, 2024+).
\item \textbf{Hueco sistemático: la vuelta 1} no tiene tiempo en FastF1 para los sprints (sí en las carreras largas). Se reconstruye con exactitud (verificada contra la FIA) desde el total oficial de F1DB, o con un desfase de salida constante por sesión para el resto de pilotos.
\item \textbf{OpenF1} (solo 2023+) coincide al milisegundo con FastF1 y rellena la v1 y las vueltas bajo SC sin tiempo, pero no es validación independiente (misma fuente Live Timing). \textbf{Jolpica} da la vuelta rápida por piloto de los 29 sprints: es la evidencia independiente (F1DB no tiene ni vuelta rápida ni paradas de sprint). formula1db.com no sirve (el TFG solo extrajo carreras largas y el sitio bloquea).
\item \textbf{Defecto actual} (independiente de las vueltas): \code{dim\_\allowbreak{}race.\allowbreak{}has\_\allowbreak{}sprint} es falso en 2021-2023 porque F1DB no trae fechas de sesión; la web oculta 12 sprints ya cargados (6 de 2021-2022 y 6 de 2023) en el selector.
\item Diseño recomendado: tipo bronze \code{sprint\_\allowbreak{}laps}, \code{int\_\allowbreak{}sprint\_\allowbreak{}laptimes} separado de \code{int\_\allowbreak{}laptimes} y \code{session\_\allowbreak{}type} en \code{fact\_\allowbreak{}laptimes}/\code{fact\_\allowbreak{}pit\_\allowbreak{}lane\_\allowbreak{}passes}; parámetro \code{?session=\allowbreak{}sprint} en \code{/\allowbreak{}laps}, \code{/\allowbreak{}stints}, \code{/\allowbreak{}pit-\allowbreak{}lane-\allowbreak{}passes}, y \code{Session\allowbreak{}Switch} en las pestañas vuelta a vuelta, neumáticos y ritmo. Esfuerzo: 6-8,5 días (mínimo viable \textasciitilde{}4).
\end{enumerate}

\section{Situación actual (revisión de código)}

\begin{itemize}
\item F1DB ya aporta resultados de sprint (\code{SPRINT\_\allowbreak{}RACE\_\allowbreak{}RESULT}) y de clasificación sprint (\code{SPRINT\_\allowbreak{}QUALIFYING\_\allowbreak{}RESULT}): \code{fact\_\allowbreak{}race\_\allowbreak{}result.\allowbreak{}session\_\allowbreak{}type in (\allowbreak{}'RACE',\allowbreak{}'SPRINT')}, \code{fact\_\allowbreak{}qualifying\_\allowbreak{}result} con \code{SPRINT\_\allowbreak{}QUALIFYING}, \code{dim\_\allowbreak{}race.\allowbreak{}has\_\allowbreak{}sprint}, \code{dim\_\allowbreak{}race.\allowbreak{}sprint\_\allowbreak{}qualifying\_\allowbreak{}format}, \code{agg\_\allowbreak{}driver\_\allowbreak{}career.\allowbreak{}sprint\_\allowbreak{}wins}.
\item API: \code{/\allowbreak{}races/\allowbreak{}\{id\}/\allowbreak{}results?session=\allowbreak{}race|sprint} y \code{/\allowbreak{}races/\allowbreak{}\{id\}/\allowbreak{}qualifying?session=\allowbreak{}qualifying|sprint\_\allowbreak{}qualifying}. Web: \code{session-\allowbreak{}switch.\allowbreak{}tsx} en resultados y clasificación.
\item Vueltas: \code{ingestion/\allowbreak{}fastf1\_\allowbreak{}loader.\allowbreak{}py} solo carga \code{get\_\allowbreak{}session(\allowbreak{}"R")} y \code{"Q"}; ningún sprint en bronze ni en la caché \code{data/\allowbreak{}cache/\allowbreak{}fastf1} (0 carpetas de sesión Sprint).
\item \code{ergast\_\allowbreak{}loader.\allowbreak{}py} usa el volcado CSV de oct-2022 (incluye \code{sprint\_\allowbreak{}results.\allowbreak{}csv}, sin vueltas de sprint); no se ingiere \code{sprint\_\allowbreak{}results}.
\item \code{int\_\allowbreak{}laptimes}, \code{int\_\allowbreak{}lap\_\allowbreak{}completeness}, \code{fact\_\allowbreak{}laptimes}, \code{fact\_\allowbreak{}pit\_\allowbreak{}lane\_\allowbreak{}passes} filtran \code{session\_\allowbreak{}type =\allowbreak{} 'RACE\_\allowbreak{}RESULT'} y claves (race\_id, driver, lap): sin dimensión de sesión.
\end{itemize}

\section{F1DB (release v2026.15.1) — medido en bronze}

\begin{itemize}
\item 29 fines de semana con \code{SPRINT\_\allowbreak{}RACE\_\allowbreak{}RESULT} (2021: 3, 2022: 3, 2023: 6, 2024: 6, 2025: 6, 2026: 5 disputados + Singapur 10-oct pendiente, ya en el calendario sin resultados).
\item \code{SPRINT\_\allowbreak{}QUALIFYING\_\allowbreak{}RESULT} solo 2023+ (23 fines de semana): en 2021-2022 la parrilla del sprint salía de la clasificación del viernes (\code{QUALIFYING\_\allowbreak{}RESULT}).
\item \code{SPRINT\_\allowbreak{}STARTING\_\allowbreak{}GRID\_\allowbreak{}POSITION} 29. NO hay \code{FASTEST\_\allowbreak{}LAP} ni \code{PIT\_\allowbreak{}STOP} de sprint: la validación de numeración (\code{int\_\allowbreak{}lap\_\allowbreak{}numbering}) y el tipado de pasos por boxes no tienen referencia F1DB.
\item \code{race\_\allowbreak{}laps} del sprint = vueltas oficiales por piloto (referencia de completitud).
\item \textbf{Defecto detectado}: \code{stg\_\allowbreak{}f1db\_\allowbreak{}\_\allowbreak{}races.\allowbreak{}has\_\allowbreak{}sprint =\allowbreak{} sprint\_\allowbreak{}race\_\allowbreak{}date is not null}, pero F1DB no trae fechas de sesión en 2021-2023 → \code{dim\_\allowbreak{}race.\allowbreak{}has\_\allowbreak{}sprint} = 0 en 2021-2023 (12 sprints con resultados en \code{fact\_\allowbreak{}race\_\allowbreak{}result} que la web no ofrece en el selector). Arreglo: derivar \code{has\_\allowbreak{}sprint} de la existencia de \code{SPRINT\_\allowbreak{}RACE\_\allowbreak{}RESULT}.
\end{itemize}

\section{FastF1 3.8.3 — prueba real (caché propia en el scratchpad)}

Calendario (\code{get\_\allowbreak{}event\_\allowbreak{}schedule}, \code{Event\allowbreak{}Format}): 30 fines de semana con sprint 2021-2026.

\begin{itemize}
\item 2021-2022 \code{sprint}: sesiones FP1 | Qualifying | FP2 | \textbf{Sprint} | Race (FastF1 normaliza el nombre oficial de 2021 «Sprint Qualifying» a \code{Sprint}).
\item 2023 \code{sprint\_\allowbreak{}shootout}: FP1 | Qualifying | \textbf{Sprint Shootout} | Sprint | Race.
\item 2024+ \code{sprint\_\allowbreak{}qualifying}: FP1 | \textbf{Sprint Qualifying} | Sprint | Qualifying | Race.
\item \code{get\_\allowbreak{}session(\allowbreak{}año,\allowbreak{} ronda,\allowbreak{} 'Sprint')} funciona para todos; la carrera sigue en \code{Session5}.
\end{itemize}

Primeras mediciones (vueltas FastF1 vs \code{race\_\allowbreak{}laps} de \code{SPRINT\_\allowbreak{}RACE\_\allowbreak{}RESULT}):

\subsection{Completitud FastF1 (29 sprints disputados, 2021-R10 → 2026-R12; medido)}

{\scriptsize\sloppy\setlength{\tabcolsep}{4pt}\renewcommand{\arraystretch}{1.15}
\begin{longtable}{@{}>{\raggedright\arraybackslash}p{0.105\dimexpr(\linewidth-64pt)\relax}>{\raggedright\arraybackslash}p{0.075\dimexpr(\linewidth-64pt)\relax}>{\raggedright\arraybackslash}p{0.085\dimexpr(\linewidth-64pt)\relax}>{\raggedright\arraybackslash}p{0.095\dimexpr(\linewidth-64pt)\relax}>{\raggedright\arraybackslash}p{0.251\dimexpr(\linewidth-64pt)\relax}>{\raggedright\arraybackslash}p{0.176\dimexpr(\linewidth-64pt)\relax}>{\raggedright\arraybackslash}p{0.172\dimexpr(\linewidth-64pt)\relax}>{\raggedright\arraybackslash}p{0.041\dimexpr(\linewidth-64pt)\relax}@{}}
\toprule
\textbf{Sprint} & \textbf{F1DB vueltas} & \textbf{FastF1 vueltas} & \textbf{Pilotos OK/total} & \textbf{Vuelta abandono (+1)} & \textbf{Tiempos nulos (v1 / resto)} & \textbf{SC / VSC / roja (vueltas-piloto)} & \textbf{Pit-in} \\
\midrule\endhead
\bottomrule\endlastfoot
2021 Gran Bretaña & 339 & 339 & 20/20 & 0 & 20 / 0 & 0/0/0 & 1 \\
2021 Italia & 342 & 343 & 19/20 & 1 & 20 / 21 & 58/0/0 & 1 \\
2021 São Paulo & 480 & 480 & 20/20 & 0 & 20 / 0 & 0/0/0 & 0 \\
2022 Emilia-Romaña & 399 & 400 & 19/20 & 1 & 20 / 0 & 77/0/0 & 1 \\
2022 Austria & 435 & 436 & 19/20 & 0 (Alonso DNS con 1 vuelta en FastF1) & 20 / 0 & 0/0/0 & 1 \\
2022 São Paulo & 468 & 469 & 19/20 & 1 & 20 / 1 & 0/0/0 & 1 \\
2023 Azerbaiyán & 308 & 308 & 20/20 & 0 & 19 / 54 & 55/38/0 & 5 \\
2023 Austria & 480 & 480 & 20/20 & 0 & 20 / 0 & 0/0/0 & 11 \\
2023 Bélgica & 208 & 209 & 19/20 & 1 & 20 / 39 & 58/0/0 & 11 \\
2023 Qatar & 318 & 322 & 16/20 & 4 & 20 / 16 & 175/0/0 & 4 \\
2023 EE. UU. & 377 & 377 & 20/20 & 0 & 20 / 0 & 0/0/0 & 1 \\
2023 São Paulo & 480 & 480 & 20/20 & 0 & 20 / 0 & 0/0/0 & 0 \\
2024 China & 378 & 378 & 20/20 & 0 & 20 / 0 & 0/0/0 & 2 \\
2024 Miami & 343 & 344 & 19/20 & 1 & 20 / 23 & 56/0/0 & 38 \\
2024 Austria & 460 & 460 & 20/20 & 0 & 20 / 0 & 0/0/0 & 0 \\
2024 EE. UU. & 380 & 380 & 20/20 & 0 & 20 / 0 & 0/0/0 & 0 \\
2024 São Paulo & 475 & 476 & 19/20 & 1 & 20 / 1 & 0/46/0 & 0 \\
2024 Qatar & 380 & 380 & 20/20 & 0 & 20 / 0 (12 vueltas sin compuesto) & 0/0/0 & 2 \\
2025 China & 380 & 380 & 20/20 & 0 & 20 / 0 & 0/0/0 & 1 \\
2025 Miami & 331 & 331 & 20/20 & 0 & 19 / 32 & 84/0/0 & 21 \\
2025 Bélgica & 297 & 297 & 20/20 & 0 & 20 / 0 & 0/0/0 & 1 \\
2025 EE. UU. & 315 & 320 & 15/20 & 5 & 20 / 40 & 150/0/0 & 3 \\
2025 São Paulo & 441 & 444 & 17/20 & 3 & 20 / 23 & 37/0/19 & 20 \\
2025 Qatar & 380 & 380 & 20/20 & 0 & 20 / 0 & 0/0/0 & 3 \\
2026 China & 396 & 397 & 21/22 & 1 & 22 / 13 & 78/0/0 & 14 \\
2026 Miami & 361 & 380 & 21/22 & 0 (Bortoleto DSQ: 19 vueltas, F1DB \code{race\_\allowbreak{}laps} nulo) & 20 / 0 & 0/0/0 & 1 \\
2026 Canadá & 490 & 491 & 21/22 & 1 & 22 / 2 & 0/0/0 & 7 \\
2026 Gran Bretaña & 372 & 372 & 22/22 & 0 & 22 / 0 & 0/0/0 & 2 \\
2026 Países Bajos & 506 & 506 & 22/22 & 0 & 22 / 1 & 0/0/0 & 2 \\
\textbf{Total} & \textbf{11 319} & \textbf{11 359} & \textbf{568/590 exactos} & 22 &  &  &  \\
\end{longtable}}

Lectura:

\begin{itemize}
\item \textbf{Cobertura 100 \%}: en los 29 sprints cada piloto tiene todas sus vueltas oficiales, sin huecos (\code{count =\allowbreak{} max(\allowbreak{}lap)}). Las 22 diferencias son la vuelta de abandono (+1, mismo criterio que \code{int\_\allowbreak{}lap\_\allowbreak{}completeness}), un DNS con vuelta registrada y un DSQ sin \code{race\_\allowbreak{}laps} en F1DB.
\item \textbf{Vuelta 1 sin tiempo SIEMPRE} en sprint (FastF1 no la calcula en sprints; en carreras largas 2023-26 solo 31 de \textasciitilde{}95 000 vueltas tienen la v1 nula). Tampoco sectores de la v1.
\item Tiempos nulos adicionales bajo SC/VSC/bandera roja (igual que en carreras largas: 1 013 de las 1 363 nulas de 2023-26 son con SC).
\item Compuestos: nombres relativos (SOFT/MEDIUM/HARD/INTERMEDIATE/WET); 1 sprint con 12 vueltas sin compuesto (2024 Qatar). Posición nula solo en la vuelta de abandono. \code{Track\allowbreak{}Status} completo.
\item F1DB no trae paradas del sprint, pero FastF1 sí registra entradas a boxes (p. ej. 2023 Austria: 11 cambios de intermedios a secos; 2024 Miami: 38 pasos bajo SC).
\end{itemize}

\subsection{Vuelta 1 del sprint: reconstruible con exactitud}

FastF1 no calcula la v1 en sprints, pero conserva \code{Lap\allowbreak{}Start\allowbreak{}Time} y \code{Time} de la v1. Medido:

\begin{itemize}
\item \code{race\_\allowbreak{}time\_\allowbreak{}millis} (F1DB) − Σ vueltas 2..n (FastF1) reproduce \textbf{exactamente} la v1 oficial de la FIA \emph{Sprint History Chart} (Imola 2022: LEC 1:31.400, VER 1:34.500, NOR 1:38.060).
\item \code{(\allowbreak{}Time − Lap\allowbreak{}Start\allowbreak{}Time)} de la v1 excede la oficial en un \textbf{desfase constante dentro de cada sesión} (160-286 ms según el sprint). Restando \code{race\_\allowbreak{}time\_\allowbreak{}penalty\_\allowbreak{}millis} (las penalizaciones de 5/10 s van incluidas en \code{race\_\allowbreak{}time\_\allowbreak{}millis}), el rango intra-sesión es 0 ms en 22 de los 24 sprints calibrables; 38 ms en 2026 China y 5 000 ms en 2026 Miami (una penalización de 5 s que F1DB no refleja en \code{race\_\allowbreak{}time\_\allowbreak{}penalty\_\allowbreak{}millis}: la mediana lo absorbe, y el caso sirve de alerta). Así se reconstruye la v1 de todos los pilotos, incluidos doblados y abandonos. En 5 sprints con SC y tiempos nulos (2021 Italia, 2023 Bakú y Bélgica, 2025 Miami y São Paulo) no hay pilotos calibradores: usar la v1 de OpenF1/FIA o el desfase típico (\textasciitilde{}260 ms) marcando \code{timing\_\allowbreak{}convention}.
\end{itemize}

\section{OpenF1 (api.openf1.org, sin clave) — prueba real}

\begin{itemize}
\item Sesiones: \code{session\_\allowbreak{}name} = \code{Sprint} (type \code{Race}) y \code{Sprint Qualifying} (type \code{Qualifying}; normaliza el «Sprint Shootout» de 2023). \textbf{Solo 2023+}: 6 sprints/año 2023-2026 (24 sesiones, Singapur 2026 incluida con \code{session\_\allowbreak{}key} 11383 aún sin datos). No cubre 2021-2022.
\item Endpoints probados por sesión: \code{laps}, \code{stints}, \code{pit}, \code{race\_\allowbreak{}control}, \code{position}, \code{session\_\allowbreak{}result} (todos responden; \textasciitilde{}6 peticiones/sprint, limitar a \textasciitilde{}2 req/s).
\end{itemize}

Validación cruzada FastF1 ↔ OpenF1 (7 sprints: 2023 Austria y Qatar, 2024 China, 2025 Bélgica y EE. UU., 2026 China y Gran Bretaña):

{\scriptsize\sloppy\setlength{\tabcolsep}{4pt}\renewcommand{\arraystretch}{1.15}
\begin{longtable}{@{}>{\raggedright\arraybackslash}p{0.127\dimexpr(\linewidth-64pt)\relax}>{\raggedright\arraybackslash}p{0.107\dimexpr(\linewidth-64pt)\relax}>{\raggedright\arraybackslash}p{0.087\dimexpr(\linewidth-64pt)\relax}>{\raggedright\arraybackslash}p{0.139\dimexpr(\linewidth-64pt)\relax}>{\raggedright\arraybackslash}p{0.101\dimexpr(\linewidth-64pt)\relax}>{\raggedright\arraybackslash}p{0.074\dimexpr(\linewidth-64pt)\relax}>{\raggedright\arraybackslash}p{0.251\dimexpr(\linewidth-64pt)\relax}>{\raggedright\arraybackslash}p{0.114\dimexpr(\linewidth-64pt)\relax}@{}}
\toprule
\textbf{Sprint} & \textbf{Vueltas FastF1} & \textbf{Emparejadas} & \textbf{Con tiempo en ambas} & \textbf{Iguales ±1 ms} & \textbf{Distintas} & \textbf{Solo OpenF1 tiene tiempo (de ellas v1)} & \textbf{Compuesto igual*} \\
\midrule\endhead
\bottomrule\endlastfoot
2023 Austria & 480 & 480 & 460 & 460 & 0 & 20 (20) & 480/480 \\
2023 Qatar & 322 & 322 & 286 & 286 & 0 & 32 (19) & 318/318 \\
2024 China & 378 & 378 & 358 & 358 & 0 & 20 (20) & 378/378 \\
2025 Bélgica & 297 & 297 & 277 & 277 & 0 & 19 (19) & 297/297 \\
2025 EE. UU. & 320 & 320 & 260 & 260 & 0 & 55 (17) & 317/318 \\
2026 China & 397 & 397 & 362 & 362 & 0 & 34 (22) & 240/240 \\
2026 Gran Bretaña & 372 & 372 & 350 & 350 & 0 & 22 (22) & 356/356 \\
\end{longtable}}

\textbackslash{}* sobre las vueltas cubiertas por algún stint de OpenF1: \textbf{0 contradicciones de compuesto}, pero OpenF1 omite stints (2026 China: le faltan los primeros stints de 157 vueltas) y numera algún stint distinto (2025 Bélgica, 11 vueltas).

Conclusiones OpenF1: mismos datos de cronometraje que FastF1 (100 \% idénticos cuando ambos tienen tiempo) → sirve para \textbf{rellenar} la v1 y las vueltas bajo SC que FastF1 deja sin tiempo, no como validación independiente (misma fuente F1 Live Timing). Su v1 es \textasciitilde{}0,35-0,44 s más larga que la oficial (mide desde otra referencia): marcarla \code{timing\_\allowbreak{}convention}, preferir la reconstrucción.

\section{Jolpica-F1 / Ergast}

\begin{itemize}
\item \code{https:\allowbreak{}/\allowbreak{}/\allowbreak{}api.\allowbreak{}jolpi.\allowbreak{}ca/\allowbreak{}ergast/\allowbreak{}f1/\allowbreak{}\{año\}/\allowbreak{}sprint/\allowbreak{}} → resultados de los 29 sprints (2021-2026), y \textbf{por piloto \code{Fastest\allowbreak{}Lap} (vuelta y tiempo)}. No existen \code{/\allowbreak{}\{año\}/\allowbreak{}\{ronda\}/\allowbreak{}sprint/\allowbreak{}laps} ni \code{/\allowbreak{}sprint/\allowbreak{}pitstops} (404); \code{/\allowbreak{}\{año\}/\allowbreak{}\{ronda\}/\allowbreak{}laps} son solo de la carrera larga.
\item Volcado CSV Ergast (oct-2022) del repo: \code{sprint\_\allowbreak{}results.\allowbreak{}csv} (100 filas = 5 sprints 2021-22) con \code{fastest\allowbreak{}Lap}/\code{fastest\allowbreak{}Lap\allowbreak{}Time}; sin vueltas de sprint.
\item Uso propuesto: la vuelta rápida por piloto de Jolpica es la \textbf{evidencia independiente} para validar tiempo y numeración (sustituye a \code{FASTEST\_\allowbreak{}LAP} de F1DB, que no existe para sprints).
\end{itemize}

\section{FIA, formula1.com y formula1db.com}

\begin{itemize}
\item FIA publica por sprint PDFs con patrón \code{fia.\allowbreak{}com/\allowbreak{}sites/\allowbreak{}default/\allowbreak{}files/\allowbreak{}\{año\}\_\allowbreak{}\{ronda\}\_\allowbreak{}\{gp\}\_\allowbreak{}f1\_\allowbreak{}s0\_\allowbreak{}timing\_\allowbreak{}sprint\{historychart|lapchart|.\allowbreak{}.\allowbreak{}.\allowbreak{}\}\_\allowbreak{}v01.\allowbreak{}pdf} (comprobados: 2022\_04\_ita \emph{Sprint History Chart}, 5 páginas con vuelta, piloto, gap, tiempo y PIT; 2026\_05\_can \emph{Sprint Lap Chart}). Es la \textbf{referencia oficial} (incluye la v1), pero son PDF de maquetación variable y con aviso de copyright restrictivo: útil para verificar casos puntuales (seed de correcciones con evidencia), no como fuente masiva.
\item formula1.com: resultados, parrilla y clasificación sprint; sin vuelta a vuelta del sprint.
\item formula1db.com: el scraping del TFG solo recorrió \code{/\allowbreak{}results/\allowbreak{}race} (no hay vueltas de sprint en \code{f1\_\allowbreak{}lap\_\allowbreak{}times.\allowbreak{}csv}); la web responde 403 y robots.txt lo prohíbe → \textbf{no es fuente para sprints}.
\end{itemize}

\section{Clasificación sprint (Shootout 2023 / Sprint Qualifying 2024+)}

\begin{itemize}
\item FastF1: \code{get\_\allowbreak{}session(\allowbreak{}a,\allowbreak{} r,\allowbreak{} 'Sprint Shootout')} (2023) y \code{'Sprint Qualifying'} (2024+) cargan vueltas (2023 Austria 355, 2024 China 218, 2025 Bélgica 185), pero \code{results} llega \textbf{sin posición ni SQ1/SQ2/SQ3} (FastF1 los toma de Ergast/Jolpica, que no los tiene) → los resultados siguen viniendo de F1DB (\code{SPRINT\_\allowbreak{}QUALIFYING\_\allowbreak{}RESULT}, ya en \code{fact\_\allowbreak{}qualifying\_\allowbreak{}result}).
\item \textbf{Trampa de alias}: \code{'SQ'} en 2021-2022 devuelve la \emph{carrera} sprint (se llamaba «Sprint Qualifying») y en 2023 falla (allí es \code{'SS'}). Usar el nombre completo según \code{Event\allowbreak{}Format}.
\item Aporte útil: telemetría de la vuelta rápida de SQ (como \code{quali\_\allowbreak{}telemetry}), no vueltas de carrera.
\end{itemize}

\section{Catálogo de sprints y cobertura esperada por fuente}

Formatos: \textbf{A} = 2021-2022 (la clasificación del viernes fija la parrilla del sprint; el sprint se llamó «Sprint Qualifying» en 2021 y «Sprint» en 2022); \textbf{B} = 2023 (Qualifying del viernes para el GP + \emph{Sprint Shootout} el sábado); \textbf{C} = 2024+ (\emph{Sprint Qualifying} el viernes; sprint y Qualifying el sábado). En 2021-2026 no se canceló ningún sprint ya disputable (Imola 2023, que iba a tener sprint, se canceló entero por las inundaciones y no figura en F1DB).

{\scriptsize\sloppy\setlength{\tabcolsep}{4pt}\renewcommand{\arraystretch}{1.15}
\begin{longtable}{@{}>{\raggedright\arraybackslash}p{0.033\dimexpr(\linewidth-72pt)\relax}>{\raggedright\arraybackslash}p{0.325\dimexpr(\linewidth-72pt)\relax}>{\raggedright\arraybackslash}p{0.033\dimexpr(\linewidth-72pt)\relax}>{\raggedright\arraybackslash}p{0.065\dimexpr(\linewidth-72pt)\relax}>{\raggedright\arraybackslash}p{0.094\dimexpr(\linewidth-72pt)\relax}>{\raggedright\arraybackslash}p{0.089\dimexpr(\linewidth-72pt)\relax}>{\raggedright\arraybackslash}p{0.111\dimexpr(\linewidth-72pt)\relax}>{\raggedright\arraybackslash}p{0.117\dimexpr(\linewidth-72pt)\relax}>{\raggedright\arraybackslash}p{0.133\dimexpr(\linewidth-72pt)\relax}@{}}
\toprule
\textbf{Año} & \textbf{Ronda / GP (F1DB \code{race\_\allowbreak{}id})} & \textbf{Fmt} & \textbf{F1DB res.} & \textbf{FastF1 vueltas} & \textbf{OpenF1} & \textbf{Jolpica res. + VR} & \textbf{Ergast CSV} & \textbf{FIA PDF} \\
\midrule\endhead
\bottomrule\endlastfoot
2021 & 10 Gran Bretaña (1045), 14 Italia (1049), 19 São Paulo (1054) & A & \si{} & \si{} medido & \no{} & \si{} & \si{} & \si{} \\
2022 & 4 Emilia-Romaña (1061), 11 Austria (1068), 21 São Paulo (1078) & A & \si{} & \si{} medido & \no{} & \si{} & \si{} (Imola, Austria) & \si{} (Imola comprobado) \\
2023 & 4 Azerbaiyán (1083), 9 Austria (1088), 12 Bélgica (1091), 17 Qatar (1096), 18 EE. UU. (1097), 20 São Paulo (1099) & B & \si{} + SQ & \si{} medido & \si{} & \si{} & \no{} & \si{} \\
2024 & 5 China (1106), 6 Miami (1107), 11 Austria (1112), 19 EE. UU. (1120), 21 São Paulo (1122), 23 Qatar (1124) & C & \si{} + SQ & \si{} medido & \si{} & \si{} & \no{} & \si{} \\
2025 & 2 China (1127), 6 Miami (1131), 13 Bélgica (1138), 19 EE. UU. (1144), 21 São Paulo (1146), 23 Qatar (1148) & C & \si{} + SQ & \si{} medido & \si{} & \si{} & \no{} & \si{} \\
2026 & 2 China (1151), 4 Miami (1153), 5 Canadá (1154), 9 Gran Bretaña (1158), 12 Países Bajos (1161) & C & \si{} + SQ & \si{} medido & \si{} & \si{} & \no{} & \si{} (Canadá comprobado) \\
2026 & 17 Singapur (1166), 10-oct & C & pendiente & pendiente & sesión creada & pendiente & \no{} & pendiente \\
\end{longtable}}

Totales: 29 sprints disputados = \textbf{11 319 vueltas oficiales} (F1DB) / 11 359 registros FastF1 (≈ +1 \% de filas sobre \code{fact\_\allowbreak{}laptimes}). OpenF1 medido en 7 de sus 23 sprints disputados; Jolpica, resultados de los 29.

Problemas identificados:

\begin{enumerate}
\item Nombres de sesión cambiantes: resolver siempre por \code{Event\allowbreak{}Format} + nombre completo, nunca por alias (\code{SQ} es ambiguo entre años).
\item v1 sin tiempo en FastF1 (reconstruible) y tiempos nulos bajo SC (rellenables con OpenF1 2023+).
\item F1DB no tiene paradas ni vuelta rápida del sprint → \code{int\_\allowbreak{}lap\_\allowbreak{}numbering}, \code{fact\_\allowbreak{}pit\_\allowbreak{}stops}, \code{fact\_\allowbreak{}pit\_\allowbreak{}lane\_\allowbreak{}passes} (\code{pass\_\allowbreak{}type=\allowbreak{}'pit\_\allowbreak{}stop'}) y \code{qa\_\allowbreak{}fastest\_\allowbreak{}laps} no tienen referencia; usar la VR de Jolpica y tipificar los pasos por boxes por SC/retirada/otros.
\item \code{dim\_\allowbreak{}race.\allowbreak{}has\_\allowbreak{}sprint} falso en 2021-2023 (defecto actual, independiente de las vueltas).
\item DSQ con \code{race\_\allowbreak{}laps} nulo (Bortoleto, 2026 Miami) y DNS con 1 vuelta en FastF1 (Alonso, 2022 Austria): conciliarlos como hace \code{int\_\allowbreak{}lap\_\allowbreak{}completeness} (\code{disqualified}, DNS excluidos).
\item Compuesto Pirelli (C1-C6): \code{int\_\allowbreak{}tyre\_\allowbreak{}compound\_\allowbreak{}mapping} depende de formula1db (carreras largas hasta 2023); para el sprint se reutiliza el mapeo de la carrera del mismo fin de semana (misma asignación de neumáticos). En 2024+ tampoco existe para la carrera: sin cambio.
\item Carga incremental: \code{completed\_\allowbreak{}races} mira \code{Session5Date\allowbreak{}Utc} (carrera); el sprint está disponible un día antes (Singapur 2026 el 10-oct). Usar la fecha de la sesión «Sprint».
\end{enumerate}

\section{Diseño propuesto}

\subsection{Ingesta}

\begin{itemize}
\item \code{fastf1\_\allowbreak{}loader.\allowbreak{}py}: nuevo tipo \code{sprint\_\allowbreak{}laps} (\code{bronze/\allowbreak{}fastf1/\allowbreak{}sprint\_\allowbreak{}laps/\allowbreak{}season=\allowbreak{}YYYY/\allowbreak{}round=\allowbreak{}RR.\allowbreak{}parquet}, con columna \code{session=\allowbreak{}'SPRINT'}) para los eventos con \code{'sprint' in Event\allowbreak{}Format}, sesión \code{'Sprint'}. Se prefiere un tipo propio a \code{laps/\allowbreak{}season=\allowbreak{}YYYY/\allowbreak{}session=\allowbreak{}S/\allowbreak{}…} porque la fuente dbt actual lee \code{fastf1/\allowbreak{}laps/\allowbreak{}*/\allowbreak{}*.\allowbreak{}parquet}: hoy no lo recogería, pero un cambio futuro del patrón a \code{**} mezclaría sprint y carrera sin avisar. Opcional: \code{sprint\_\allowbreak{}quali\_\allowbreak{}telemetry} (sesiones «Sprint Shootout»/«Sprint Qualifying», mismo patrón que \code{quali\_\allowbreak{}telemetry}).
\item Nuevo \code{openf1\_\allowbreak{}loader.\allowbreak{}py} (2023+): \code{laps}, \code{stints}, \code{pit}, \code{race\_\allowbreak{}control} de las sesiones \code{Sprint} → \code{bronze/\allowbreak{}openf1/\allowbreak{}\{tipo\}/\allowbreak{}season=\allowbreak{}YYYY/\allowbreak{}round=\allowbreak{}RR.\allowbreak{}parquet}; \code{session\_\allowbreak{}key} resuelto con \code{/\allowbreak{}sessions?year=\allowbreak{}…} (filtrar \code{session\_\allowbreak{}name =\allowbreak{}=\allowbreak{} 'Sprint'}) y casado con F1DB por fecha y país.
\item Jolpica: \code{jolpica\_\allowbreak{}loader.\allowbreak{}py} con \code{/\allowbreak{}\{año\}/\allowbreak{}sprint/\allowbreak{}} (resultados + \code{Fastest\allowbreak{}Lap}), paginando (\code{limit=\allowbreak{}100}, \code{offset}). Sirve además para contrastar los resultados de sprint de F1DB.
\end{itemize}

\subsection{dbt}

\begin{itemize}
\item Staging: \code{stg\_\allowbreak{}fastf1\_\allowbreak{}\_\allowbreak{}sprint\_\allowbreak{}laps} (igual que \code{stg\_\allowbreak{}fastf1\_\allowbreak{}\_\allowbreak{}laps} + \code{lap\_\allowbreak{}start\_\allowbreak{}time\_\allowbreak{}ms}), \code{stg\_\allowbreak{}openf1\_\allowbreak{}\_\allowbreak{}laps}, \code{stg\_\allowbreak{}jolpica\_\allowbreak{}\_\allowbreak{}sprint\_\allowbreak{}results}.
\item \code{int\_\allowbreak{}fastf1\_\allowbreak{}sprint\_\allowbreak{}laps} (prototipo en el scratchpad): claves F1DB + reconstrucción de la v1.
\item \code{int\_\allowbreak{}sprint\_\allowbreak{}laptimes}: FastF1 principal; OpenF1 rellena tiempos nulos (\code{source=\allowbreak{}'openf1'}); \code{validation\_\allowbreak{}status} con la misma semántica que \code{int\_\allowbreak{}laptimes}, más un valor nuevo:
\begin{itemize}
\item \code{confirmed}: FastF1 = OpenF1 ±1 ms (documentar que comparten la fuente primaria) o = VR de Jolpica;
\item \code{derived} (\textbf{nuevo}): v1 reconstruida desde el total oficial de F1DB;
\item \code{timing\_\allowbreak{}convention}: v1 de OpenF1 o por desfase de sesión;
\item \code{disputed}: la VR de Jolpica no coincide con ninguna vuelta del piloto;
\item \code{single\_\allowbreak{}source}: resto. Numeración: evidencia = VR de Jolpica (±1 ms) en lugar de \code{FASTEST\_\allowbreak{}LAP} de F1DB.
\end{itemize}
\item \textbf{Hecho}: añadir \code{session\_\allowbreak{}type} (\code{RACE}/\code{SPRINT}) a \code{fact\_\allowbreak{}laptimes} (union de \code{int\_\allowbreak{}laptimes}, sin tocar, + \code{int\_\allowbreak{}sprint\_\allowbreak{}laptimes}) y a \code{fact\_\allowbreak{}pit\_\allowbreak{}lane\_\allowbreak{}passes}; clave única \code{(\allowbreak{}race\_\allowbreak{}id,\allowbreak{} session\_\allowbreak{}type,\allowbreak{} driver\_\allowbreak{}id,\allowbreak{} driver\_\allowbreak{}number,\allowbreak{} lap\_\allowbreak{}number)}. Se elige columna en vez de hecho separado porque la API y la web ya tratan la sesión como parámetro (\code{fact\_\allowbreak{}race\_\allowbreak{}result} hace lo mismo) y las consultas de vueltas, stints y ritmo se reutilizan tal cual. Contrapartida: todo consumidor debe filtrar la sesión: 3 consultas de \code{api/\allowbreak{}app/\allowbreak{}routers/\allowbreak{}races.\allowbreak{}py} (\code{laps}, \code{stints}, \code{pit-\allowbreak{}lane-\allowbreak{}passes}), \code{qa\_\allowbreak{}fastest\_\allowbreak{}laps}, \code{qa\_\allowbreak{}pit\_\allowbreak{}stops}, \code{qa\_\allowbreak{}summary} e \code{int\_\allowbreak{}lap\_\allowbreak{}completeness}. Alternativa de menor riesgo: \code{fact\_\allowbreak{}sprint\_\allowbreak{}laptimes} separado.
\item \code{int\_\allowbreak{}lap\_\allowbreak{}completeness} parametrizado por sesión (\code{SPRINT\_\allowbreak{}RACE\_\allowbreak{}RESULT}) → \code{quality\_\allowbreak{}status}.
\item \code{stg\_\allowbreak{}f1db\_\allowbreak{}\_\allowbreak{}races.\allowbreak{}has\_\allowbreak{}sprint}: derivarlo de la existencia de \code{SPRINT\_\allowbreak{}RACE\_\allowbreak{}RESULT} (arreglo inmediato, útil aunque no se haga lo demás).
\item Tests: \code{unique\_\allowbreak{}combination} con \code{session\_\allowbreak{}type}; \code{accepted\_\allowbreak{}values} de \code{session\_\allowbreak{}type}; singular \code{assert\_\allowbreak{}sprint\_\allowbreak{}laps\_\allowbreak{}match\_\allowbreak{}f1db} (máx. vuelta = \code{race\_\allowbreak{}laps} salvo vuelta de abandono/DSQ/DNS); \code{assert\_\allowbreak{}sprint\_\allowbreak{}lap\_\allowbreak{}sum\_\allowbreak{}matches\_\allowbreak{}total} (Σ vueltas = \code{race\_\allowbreak{}time − penalización} ±1 ms para los clasificados en la vuelta del líder); umbral de sprint en \code{assert\_\allowbreak{}quality\_\allowbreak{}thresholds}; \code{qa\_\allowbreak{}sprint\_\allowbreak{}fastest\_\allowbreak{}laps} (VR Jolpica frente a vueltas).
\end{itemize}

\subsection{API}

\begin{itemize}
\item \code{GET /\allowbreak{}races/\allowbreak{}\{id\}/\allowbreak{}laps?session=\allowbreak{}race|sprint}, \code{/\allowbreak{}stints?session=\allowbreak{}…}, \code{/\allowbreak{}pit-\allowbreak{}lane-\allowbreak{}passes?session=\allowbreak{}…} (\code{Literal["race",\allowbreak{}"sprint"] =\allowbreak{} "race"}, como \code{/\allowbreak{}results}); lista vacía si no hay sprint. \code{/\allowbreak{}pitstops} sigue siendo solo de carrera (F1DB no tiene paradas de sprint). Regenerar \code{web/\allowbreak{}src/\allowbreak{}lib/\allowbreak{}api/\allowbreak{}openapi.\allowbreak{}json} y \code{schema.\allowbreak{}d.\allowbreak{}ts}; fixtures de \code{api/\allowbreak{}tests} con un sprint.
\end{itemize}

\subsection{Web}

\begin{itemize}
\item Pestañas \code{lap-\allowbreak{}chart}, \code{tyres}, \code{pace} (y \code{pitstops} para los pasos por boxes): \code{Session\allowbreak{}Switch} carrera/sprint cuando \code{race.\allowbreak{}has\_\allowbreak{}sprint}, con \code{?session=\allowbreak{}sprint} (sin JS, como en resultados). \code{get\allowbreak{}Driver\allowbreak{}Names(\allowbreak{}race\allowbreak{}Id,\allowbreak{} session)} debe usar \code{get\allowbreak{}Results(\allowbreak{}race\allowbreak{}Id,\allowbreak{} "sprint")} para el orden. Textos en \code{dictionaries/\allowbreak{}\{es,\allowbreak{}en\}.\allowbreak{}json} y aviso de «vuelta 1 derivada».
\end{itemize}

\subsection{Esfuerzo estimado}

{\small\sloppy\setlength{\tabcolsep}{5pt}\renewcommand{\arraystretch}{1.15}
\begin{longtable}{@{}>{\raggedright\arraybackslash}p{0.709\dimexpr(\linewidth-20pt)\relax}>{\raggedright\arraybackslash}p{0.291\dimexpr(\linewidth-20pt)\relax}@{}}
\toprule
\textbf{Bloque} & \textbf{Días} \\
\midrule\endhead
\bottomrule\endlastfoot
Loader FastF1 sprint + backfill de 29 sprints (\textasciitilde{}1-2 min/sprint con caché fría) & 0,5-1 \\
Loaders OpenF1 (2023+) y Jolpica sprint & 1 \\
dbt: staging, \code{int\_\allowbreak{}fastf1\_\allowbreak{}sprint\_\allowbreak{}laps} (v1), \code{int\_\allowbreak{}sprint\_\allowbreak{}laptimes}, \code{session\_\allowbreak{}type} en hechos, \code{has\_\allowbreak{}sprint} & 2-3 \\
Tests dbt + QA & 1 \\
API (3 endpoints + tests + OpenAPI) & 0,5-1 \\
Web (selector en 3-4 pestañas, i18n, e2e Playwright) & 1-1,5 \\
\textbf{Total} & \textbf{6-8,5 (≈ 1,5 semanas)} \\
\end{longtable}}

Versión mínima (solo FastF1 + v1 derivada + \code{session\_\allowbreak{}type} + API/web, sin OpenF1/Jolpica): \textasciitilde{}4 días.

\section{Riesgos}

\begin{itemize}
\item \textbf{Regresión por cambio de grano} de \code{fact\_\allowbreak{}laptimes}: una consulta sin filtro de sesión mezcla sprint y carrera (vueltas 1..24 duplicadas, ritmo sesgado). Mitigar con test de grano y revisando los consumidores listados.
\item \textbf{Validación poco independiente}: FastF1 y OpenF1 salen del mismo F1 Live Timing; lo independiente es la VR de Jolpica y el total de F1DB. La FIA (PDF), solo para casos puntuales.
\item Dependencia de APIs de terceros (Live Timing vía FastF1, límites de OpenF1, Jolpica \textasciitilde{}200 peticiones/hora). Mitigan la caché y el bronze por sprint.
\item v1 no calibrable en 5 sprints con SC y tiempos nulos → \code{timing\_\allowbreak{}convention}.
\item F1DB puede no reflejar la penalización en \code{race\_\allowbreak{}time\_\allowbreak{}penalty\_\allowbreak{}millis} (2026 Miami): la v1 derivada saldría 5 s desplazada. Mitigar con la mediana del desfase y un test de rango (v1 entre 0,9× y 1,6× la v2 del piloto) que dispare el uso del desfase.
\item Licencia: los PDF de la FIA prohíben la reproducción sin permiso; no redistribuirlos.
\end{itemize}

\section{Prototipo (scratchpad \code{…/\allowbreak{}730e2b1e-\allowbreak{}…/\allowbreak{}scratchpad/\allowbreak{}sprint/\allowbreak{}})}

\begin{itemize}
\item \code{schedule.\allowbreak{}py} → \code{sprint\_\allowbreak{}schedule\_\allowbreak{}fastf1.\allowbreak{}csv} (30 fines de semana con sprint 2021-2026).
\item \code{ff1\_\allowbreak{}sprint\_\allowbreak{}probe.\allowbreak{}py} → carga \code{get\_\allowbreak{}session(\allowbreak{}a,\allowbreak{} r,\allowbreak{} 'Sprint')}, escribe \code{bronze\_\allowbreak{}proto/\allowbreak{}fastf1/\allowbreak{}laps/\allowbreak{}season=\allowbreak{}YYYY/\allowbreak{}session=\allowbreak{}S/\allowbreak{}round=\allowbreak{}RR.\allowbreak{}parquet} y mide frente a F1DB.
\item \code{measure\_\allowbreak{}all.\allowbreak{}py} → \code{coverage\_\allowbreak{}fastf1\_\allowbreak{}sprint.\allowbreak{}csv} (tabla de completitud); \code{anomalies.\allowbreak{}py}.
\item \code{openf1\_\allowbreak{}sessions.\allowbreak{}py}, \code{openf1\_\allowbreak{}probe.\allowbreak{}py} → \code{openf1\_\allowbreak{}sprint\_\allowbreak{}sessions.\allowbreak{}csv} y bronze OpenF1.
\item \code{compare\_\allowbreak{}ff1\_\allowbreak{}openf1.\allowbreak{}py} (validación cruzada), \code{jolpica\_\allowbreak{}probe.\allowbreak{}py}, \code{sq\_\allowbreak{}probe.\allowbreak{}py}, \code{lap1\_\allowbreak{}derivation.\allowbreak{}py} (calibración de la v1), \code{fia\_\allowbreak{}2022\_\allowbreak{}04\_\allowbreak{}ita\_\allowbreak{}sprinthistorychart.\allowbreak{}txt}.
\item \code{proposal/\allowbreak{}fastf1\_\allowbreak{}loader\_\allowbreak{}sprint.\allowbreak{}py} y \code{proposal/\allowbreak{}int\_\allowbreak{}fastf1\_\allowbreak{}sprint\_\allowbreak{}laps.\allowbreak{}sql} (código propuesto).
\end{itemize}

Código clave (reconstrucción de la v1, del prototipo dbt):

\begin{codigo}
session_offset as (
    select race_id, median(lap1_from_start_ms - (race_time_ms - rest_ms)) as start_offset_ms
    from per_driver                 -- race_time_ms = race_time_millis - race_time_penalty_millis
    where race_time_ms is not null and rest_laps = race_laps - 1
    group by race_id
)
-- v1 = race_time_ms - suma(v2..n)   si el piloto tiene total oficial y vueltas 2..n completas
--    = (Time - LapStartTime de la v1) - start_offset_ms   en otro caso
\end{codigo}

Carga de la sesión (loader propuesto):

\begin{codigo}
def _load_sprint(event, season):
    session = event.get_session("Sprint")  # válido en los tres formatos 2021+
    session.load(laps=True, telemetry=False, weather=False, messages=False)
    laps = _with_keys(timedeltas_to_millis(pd.DataFrame(session.laps)), event, season)
    laps.insert(3, "session", "SPRINT")
    write_parquet(laps, race_path("sprint_laps", season, int(event["RoundNumber"])))
\end{codigo}
